\begin{abstract}
\noindent\textbf{Question.} Thin composite panels harden as they vibrate: the resonance drifts
upward with amplitude and, under forcing, folds and jumps. Can the \emph{spatial distribution}
of fiber orientation reduce this hardening at fixed linear frequency and mass, by what
mechanism, and how does it compare with a uniform ply angle and with the stacking sequence?

\noindent\textbf{Hypothesis.} The frequency and the hardening integrate the same rotated
stiffness against two different strain fields of the fundamental mode, its curvature and its
in-plane-relaxed stretch. Where these fields differ in principal direction, a fiber field can
serve both pointwise; a uniform layup can only compromise.

\noindent\textbf{Analysis.} A local Fourier representation predicts no lever for isotropic
material, invariance of a single-ply beam under any uniform angle, and a pointwise first-order
optimality condition. We test the predictions on laminated von K\'arm\'an beams and plates, with
the frequency-normalized hardening $J$ of a spectral-submanifold (SSM) reduction (in-plane
relaxation exact for the von K\'arm\'an model, mode-adaptive, analytic adjoint), the frequency
as an equality constraint and the tow-curvature limit imposed pointwise.

\noindent\textbf{Results.} On a clamped beam, symmetric steering cuts the hardening by up to
$\nBInfHeldRed\%$ by placing compliant fibers at the inflection points. On a simply supported
plate, a curvature-feasible idealized tow field ($R_{\min}=\SI{20}{mm}$) brings $J$ to $\nFTen$ of
the unidirectional value at the unidirectional frequency; with the frequency re-matched in
refined and independent models the ratio is $\nFEmatchedJ$--$\nFullMatchedEPMC$. At that frequency no
other uniform symmetric laminate exists, and above it steering undercuts every uniform class
surveyed. The lever vanishes for isotropic material, grows with anisotropy, saturates once tows
turn on the strain-field length scale, and persists for clamped edges ($\nGenClJ$ against
$\nGenClUnif$ for the best uniform laminate), a square plate ($\nGenSqJ$) and in-plane movable edges
($\nGenMvJ$); it is weak on a 2:1 plate whose near-cylindrical mode leaves little conflict between the
two fields. The jump-free amplitude widens by $\nCacRatio\times$, as $a_c\propto J^{-1/2}$ predicts.

\noindent\textbf{Conclusion.} Spatial fiber steering is a mechanistically understood lever on
nonlinear stiffness at fixed linear frequency and mass, set by the material anisotropy, the conflict
between the two strain fields and the tow radius. Laying the idealized fields requires tow cut/add;
an as-built defect model keeps $J/J_{\UD}$ between $\nAsBuiltJmin$ and $\nAsBuiltJmax$, but no longer at the
held frequency and mass (frequency $\nAsBuiltFmin$ to $\nAsBuiltFmax$, mass $\nAsBuiltMmin$ to
$\nAsBuiltMmax$ of UD).
\end{abstract}

\noindent\textbf{Keywords:} variable-stiffness composites; fiber steering; geometric nonlinearity;
spectral submanifolds; backbone curve; adjoint optimization; automated fiber placement.
